% docs/manuscript/sections/framework_v2.tex
% Section III v2 -- unified formal framework (stream S1).
% Replaces the v63 sections sec:race, sec:hydra, sec:starvation, sec:decoupling.

\section{Formal Framework}\label{sec:framework}

\subsection{Stream, Labels, and Observability}\label{sec:observability}

We consider a stream of pairs $(x_t, y_t)_{t \ge 1}$ drawn from a joint law
$P_t$ that is stationary before an unknown change point $\tau^*$ and
stationary after it. A classifier $C$ emits $\hat{y}_t$ and the
\emph{prediction error} is $e_t := \mathbf{1}[\hat{y}_t \ne y_t]$.
Write $\bar{e}_t := \mathbb{E}[e_t]$ for its expectation under the joint law
of stream and classifier, $p_0 := \bar{e}_t$ for $t < \tau^*$ for the
pre-change error rate, and $\Delta e := \sup_{t \ge \tau^*} \bar{e}_t - p_0$
for the error jump.

\begin{assumption}[Observability split]\label{ass:obs}
  The change point $\tau^*$ and the post-change law $P_{\tau^*+}$ are
  unobservable. The error $e_t$ is observable exactly when $y_t$ is revealed.
\end{assumption}

Assumption~\ref{ass:obs} fixes the semantics of every subsequent statement:
what is hidden is the drift \emph{state}, never the error \emph{stream}.

\begin{definition}[Label protocol]\label{def:labels}
  A label protocol $\Pi$ reveals $y_t$ at time $t + \ell$ with probability
  $p \in (0,1]$, independently across $t$. Immediate supervision is
  $\Pi(\ell{=}0, p{=}1)$. The \emph{label clock} indexes the subsequence of
  revealed labels; the \emph{stream clock} indexes $t$.
\end{definition}

\begin{definition}[Monitor and windowed false-alarm level]\label{def:monitor}
  A monitor $D$ is a stopping time $\tau_{\mathrm{det}}$ adapted to the
  filtration generated by the observed errors. Its \emph{windowed
  false-alarm level} at horizon $W$ is
  $\alpha := \mathbb{P}_0(\tau_{\mathrm{det}} \le W)$, the probability of
  alarming within $W$ consecutive steps under the no-change law $P_0$.
\end{definition}

Definition~\ref{def:monitor} replaces family-specific tuning constants by a
single quantity shared by every sequential test. Under the geometric
approximation standard in statistical process control, $\alpha \approx W /
\mathrm{ARL}_0$, where $\mathrm{ARL}_0 := \mathbb{E}_0[\tau_{\mathrm{det}}]$.

\begin{lemma}[Label-clock invariance]\label{lem:clock}
  If $C$ and $D$ consume the same label subsequence, then the quantities
  $A$ and $R$ of Definitions~\ref{def:budget} and~\ref{def:requirement} are
  invariant under $\Pi(\ell, p)$ when evaluated in the label clock.
\end{lemma}

\begin{proof}
  A uniform latency $\ell$ shifts both the adaptation and the detection
  filtrations by the same $\ell$, leaving every difference of stopping times
  unchanged. Under thinning at rate $p$, $C$ requires the same number of
  labelled examples to complete adaptation, so the recovery window dilates by
  $1/p$ in the stream clock, while the number of monitor updates inside that
  window contracts by $p$. The two factors cancel.
\end{proof}

Lemma~\ref{lem:clock} fails only when the monitor and the classifier read
different filtrations, e.g.\ an unsupervised feature-space monitor paired
with a supervised learner.

\subsection{Adaptation Times}\label{sec:times}

Let $M$ be the ensemble size and $N_{\mathrm{swap}}(t)$ the number of base
learners replaced in $(\tau^*, t]$.

\begin{definition}[Adaptation times]\label{def:times}
  For $q \in (0,1]$ and $\rho > 0$,
  \begin{align*}
    \tau_{\mathrm{swap}}^{(q)} &:= \inf\{\, t \ge \tau^* :
      N_{\mathrm{swap}}(t)/M \ge q \,\}, \\
    \tau_{\mathrm{err}}(\rho) &:= \inf\{\, t \ge \tau^* :
      \bar{e}_s \le p_0 + \rho \ \ \forall s \ge t \,\}, \\
    \tau_{\mathrm{erase}} &:= \tau_{\mathrm{err}}(\delta_P), \qquad
    W := \tau_{\mathrm{erase}} - \tau^*,
  \end{align*}
  where $\delta_P > 0$ is the drift tolerance of the monitor and $W$ is the
  \emph{exploitable transient}.
\end{definition}

\begin{assumption}[No recovery without replacement]\label{ass:repair}
  On $[\tau^*, \tau_{\mathrm{swap}}^{(1/M)})$ the ensemble error satisfies
  $\bar{e}_t > p_0 + \delta_P$.
\end{assumption}

\begin{proposition}[Ordering]\label{prop:order}
  $\rho \mapsto \tau_{\mathrm{err}}(\rho)$ is non-increasing and
  $q \mapsto \tau_{\mathrm{swap}}^{(q)}$ non-decreasing. Under
  Assumption~\ref{ass:repair},
  $\tau_{\mathrm{swap}}^{(1/M)} \le \tau_{\mathrm{erase}}
   \le \tau_{\mathrm{swap}}^{(1)}$.
\end{proposition}

\begin{definition}[Dilation ratio]\label{def:kappa}
  $\kappa := (\tau_{\mathrm{erase}} - \tau^*) /
             (\tau_{\mathrm{swap}}^{(1/M)} - \tau^*) \ge 1$.
\end{definition}

The first swap is retained as an instrumented lower bound on $W$; $\kappa$ is
measured, never assumed.

\subsection{Proof Budget}\label{sec:budget}

\begin{definition}[Proof budget]\label{def:budget}
  \begin{equation}\label{eq:budget}
    A \;:=\; \sum_{t = \tau^*+1}^{\tau_{\mathrm{erase}}}
             \bigl(\bar{e}_t - p_0 - \delta_P\bigr)^{+} .
  \end{equation}
\end{definition}

$A$ is the integrated excess error above the monitor's tolerance over the
exploitable transient, measured in units of excess errors. It depends on the
classifier and the stream, never on the monitor beyond the scalar
$\delta_P$. A rectangular transient of height $\Delta e$ and duration $W$
gives $A = (\Delta e - \delta_P)\,W$.

\begin{proposition}[Budget invariance]\label{prop:invariance}
  If the adaptation time obeys a power law
  $W(\Delta e) = K (\Delta e)^{-\alpha}$ with $\alpha = 1$, then
  $A = K (1 - \delta_P/\Delta e) \to K$ as $\Delta e / \delta_P \to \infty$:
  the proof budget is asymptotically independent of the drift magnitude.
\end{proposition}

Proposition~\ref{prop:invariance} is the formal content of the paradox.
Increasing the magnitude of a drift does not increase the evidence delivered
to an external monitor, because adaptation accelerates in proportion.

\subsection{Proof Requirement}\label{sec:requirement}

\begin{definition}[Proof requirement]\label{def:requirement}
  For a monitor $D$, a miss level $\varepsilon \in (0,1)$ and a windowed
  false-alarm level $\alpha$,
  \[
    R(D, \varepsilon, \alpha) := \inf\bigl\{\, a > 0 :
      A \ge a \Rightarrow
      \mathbb{P}(\tau_{\mathrm{det}} \le \tau^* + W) \ge 1 - \varepsilon
      \,\bigr\},
  \]
  the infimum being taken over error trajectories of window $W$ compatible
  with budget $A$.
\end{definition}

\begin{definition}[Blind spot]\label{def:blindspot}
  A pipeline $(C, D)$ exhibits a \emph{blind spot} at $(\Delta e, W)$ when
  $A < R(D, \varepsilon, \alpha)$.
\end{definition}

Definition~\ref{def:blindspot} compares two quantities in the same unit and
contains no implementation constant. It supersedes any condition coupling an
ADWIN clock to a CUSUM threshold.

\begin{theorem}[Universal detection floor]\label{thm:floor}
  Let $\Delta_{\max} := \max_{t \in (\tau^*, \tau^*+W]}
  (\bar{e}_t - p_0)$. Every monitor satisfying
  Definition~\ref{def:monitor} at level $\alpha$ and detecting with
  probability at least $1 - \varepsilon$ within $W$ obeys
  \begin{equation}\label{eq:floor}
    A \;\ge\; \frac{p_0(1-p_0)}{\Delta_{\max}}\,
              d\bigl(1-\varepsilon \,\Vert\, \alpha\bigr) \;-\; W \delta_P ,
  \end{equation}
  where $d(a \Vert b) := a \ln(a/b) + (1-a)\ln\frac{1-a}{1-b}$ is the binary
  Kullback--Leibler divergence.
\end{theorem}

\begin{proof}
  Let $P_1$ be the post-change law and $P_0$ the no-change law restricted to
  the $W$ error observations. Set $u_t := \bar{e}_t - p_0 \ge 0$. Then
  $u_t \le (u_t - \delta_P)^{+} + \delta_P$, hence
  $\sum_t u_t \le A + W\delta_P$ and, since $u_t \le \Delta_{\max}$,
  $\sum_t u_t^2 \le \Delta_{\max}(A + W\delta_P)$. Chain rule and the bound
  $d(x \Vert y) \le (x-y)^2 / (y(1-y))$ give
  $\mathrm{KL}(P_1 \Vert P_0) \le \Delta_{\max}(A + W\delta_P) /
  (p_0(1-p_0))$. Applying the data-processing inequality to the event
  $E := \{\tau_{\mathrm{det}} \le \tau^*+W\}$, which has $P_1$-probability at
  least $1-\varepsilon$ and $P_0$-probability at most $\alpha$, yields
  $\mathrm{KL}(P_1 \Vert P_0) \ge d(1-\varepsilon \Vert \alpha)$.
  Rearranging gives~\eqref{eq:floor}.
\end{proof}

Theorem~\ref{thm:floor} separates two regimes. When $A$ falls below the
right-hand side of~\eqref{eq:floor}, no monitor of that false-alarm level
detects, and the blind spot is information-theoretic. When $A$ exceeds it
while remaining below $R(D,\varepsilon,\alpha)$, the blind spot is a property
of $D$'s calibration, and a different monitor family escapes it.

\paragraph{Family instantiations.} Write $\lambda(\alpha)$ for the CUSUM
threshold calibrated to level $\alpha$, $n_{\mathrm{stat}}$ for the KSWIN
statistic size, and $\sigma^2 \le 1/4$ for the per-step error variance.
Up to universal constants,
\begin{align}
  R_{\mathrm{CUSUM}} &= \lambda(\alpha)
     + \sqrt{\tfrac{W}{2}\ln \tfrac{1}{\varepsilon}},
     \label{eq:Rcusum} \\
  R_{\mathrm{ADWIN}} &= \sqrt{\tfrac{W}{2}\ln \tfrac{4W}{\alpha}}
     + \sqrt{\tfrac{W}{2}\ln \tfrac{1}{\varepsilon}},
     \label{eq:Radwin} \\
  R_{\mathrm{KSWIN}} &= \sqrt{n_{\mathrm{stat}} \ln \tfrac{2}{\alpha}}
     + \sqrt{\tfrac{W}{2}\ln \tfrac{1}{\varepsilon}} .
     \label{eq:Rkswin}
\end{align}

\begin{corollary}[Family split]\label{cor:split}
  The false-alarm price enters $R_{\mathrm{CUSUM}}$ additively and
  independently of $W$, and enters $R_{\mathrm{ADWIN}}$ and
  $R_{\mathrm{KSWIN}}$ inside a square root of $W$. Hence, as $W \to 0$,
  $R_{\mathrm{CUSUM}} \to \lambda(\alpha) > 0$ whereas
  $R_{\mathrm{ADWIN}} \to 0$. Cumulative monitors carry an incompressible
  requirement that short transients cannot meet; windowed monitors do not.
\end{corollary}

Corollary~\ref{cor:split} derives the empirically observed separation
between the cumulative and the windowed families from the scaling of
$R$ in $W$, rather than positing it as a taxonomy.